%%%%%%%%%%%%%%%%%%%%%%%%%%%%%%%%%%%%%%%%%%%%%%%%%%%%%%%%%%%%%%%%%%%
%%% Documento LaTeX - Comandos
%%%%%%%%%%%%%%%%%%%%%%%%%%%%%%%%%%%%%%%%%%%%%%%%%%%%%%%%%%%%%%%%%%%

% 1. Informacion del documento
\newcommand{\tfgtitlename}{Evaluación de Seguridad en Ecosistemas Widevine DRM: Diseño de una Arquitectura de Defensa en Profundidad para el \textit{Streaming} en Directo}
\newcommand{\tfgtitlenameENG}{Security Evaluation of Widevine DRM Ecosystems: Designing a Defense-in-Depth Architecture for Live Streaming}
\newcommand{\tfgauthorname}{Jesús Ramírez Gálvez}
\newcommand{\tfgtutorname}{Antonio Muñoz Gallego}
\newcommand{\tfganno}{2026}

% 2. Datos de portada en el formato del TFG original
\newcommand{\nombreautor}{\tfgauthorname}
\newcommand{\nombretutor}{\tfgtutorname}
\newcommand{\titulotrabajo}{\tfgtitlename}
\newcommand{\escuela}{Escuela Técnica Superior\\de Ingeniería Informática}
\newcommand{\escuelalargo}{Escuela Técnica Superior de Ingeniería Informática}
\newcommand{\universidad}{Universidad Rey Juan Carlos}
\newcommand{\fecha}{Fecha}
\newcommand{\grado}{Máster en Ciberseguridad}
\newcommand{\curso}{Curso 2025-2026}
\newcommand{\logoUniversidad}{logoURJC.pdf}

% 3. Comandos generales de texto
\newcommand{\R}{\textsuperscript{\textregistered}}
\newcommand{\C}{\textsuperscript{\copyright}}
\newcommand{\TM}{\texttrademark}
\newcommand{\tit}{\textit}
\newcommand{\tbf}{\textbf}
\newcommand{\ttw}[1]{\texttt{#1}}
\newcommand{\textittt}[1]{\textit{\texttt{#1}}}
\newcommand{\textittw}{\textittt}
\newcommand{\tittw}{\textittw}
\newcommand{\tbftw}[1]{\tbf{\ttw{#1}}}
\newcommand{\nli}{\\ \indent}
\newcommand{\mailto}[1]{\href{mailto:#1}{#1}}
\newcommand{\miindex}[1]{#1\index{#1}}
\newcommand{\hs}{\hspace}
\newcommand{\vs}{\vspace}

% 4. Abreviaturas matematicas
\newcommand{\realnumbers}{\mathbb R}
\newcommand{\naturalnumbers}{\mathbb N}
\newcommand{\integernumbers}{\mathbb Z}
\newcommand{\rationalnumbers}{\mathbb Q}
\newcommand{\complexnumbers}{\mathbb R}
\newcommand{\irrationalnumbers}{\mathbb I}
\newcommand{\doublebar}[1]{\bar{\bar{#1}}}
\newcommand{\transpuesta}{\mbox{\tiny $\mathsf{T}$}}

% 5. Paginas, capitulos e indices
\newcommand\blankpage{%
    \newpage
    \null
    \thispagestyle{empty}%
    \newpage}

\newcommand{\sectionx}[1]{%
    \section*{#1}
    \addcontentsline{toc}{section}{#1}}

\newcommand{\sectiony}[1]{%
    \section*{#1}}

\newcommand{\chapterx}[1]{%
    \chapter*{#1}
    \addcontentsline{toc}{chapter}{#1}}

\newcommand{\chapterbegin}[1]{%
    \pagestyle{fancy}
    \chapter{#1}}

\newcommand{\chapterbeginx}[1]{%
    \pagestyle{fancy}
    \chapterx{#1}}

\newcommand{\chapterend}{%
    \clearpage
    \thispagestyle{empty}}

% 6. Colores, enlaces y textos automaticos
\definecolor{BlueLink}{rgb}{0.165,0.322,0.745}
\definecolor{PinkLink}{rgb}{0.8,0.22,0.5}
\definecolor{gray}{rgb}{0.6,0.6,0.6}
\definecolor{codegray}{rgb}{0.95,0.95,0.95}

\hypersetup{hidelinks,pageanchor=true,colorlinks,citecolor=PinkLink,urlcolor=black,linkcolor=BlueLink}

\addto{\captionsspanish}{\renewcommand{\bibname}{Bibliografía}}
\addto\captionsspanish{%
    \def\tablename{Tabla}
    \def\listtablename{Índice de tablas}}
\addto\captionsspanish{\renewcommand{\contentsname}{Índice de contenidos}}

\floatname{algorithm}{Algoritmo}

% 7. Entornos auxiliares heredados de la plantilla base
\newcommand{\benu}{\begin{enu}}
\newcommand{\eenu}{\end{enu}}

\newcommand{\code}[3]{%
\begin{minipage}{0.96\linewidth}
\begin{center}
    \lstinputlisting[
        language   = Python,
        frame      = lines,
        framexleftmargin = 5mm,
        caption    = {#1},
        label      = {#3},
    ]{#2}
\end{center}
\end{minipage}}

% 8. Estilo de codigo del TFG original
\definecolor{bg}{rgb}{0.95,0.95,0.95}
\definecolor{mydeepteal}{rgb}{0.16,0.22,0.23}
\definecolor{myteal}{rgb}{0.31,0.44,0.46}
\definecolor{mymediumteal}{rgb}{0.41,0.58,0.60}

\DeclareFixedFont{\ttb}{T1}{txtt}{bx}{n}{12}
\DeclareFixedFont{\ttm}{T1}{txtt}{m}{n}{12}

\newcommand*{\FormatDigit}[1]{\textcolor{black}{#1}}

\newcommand\mypythonstyle{\lstset{
language=Python,
basicstyle=\ttfamily\small,
otherkeywords={self},
keywordstyle=\bfseries\ttfamily\color{myteal},
commentstyle=\itshape\color{myteal},
stringstyle=\color{mydeepteal},
emph={MyClass,__init__},
emphstyle=\ttb\color{mydeepteal},
showstringspaces=false,
backgroundcolor=\color{bg},
rulecolor=\color{bg},
breaklines=true,
numbers=left,
numbersep=5pt,
numberstyle=\tiny,
tabsize=4,
xleftmargin=1em,
frame=single,
framesep=3pt,
framextopmargin=0pt,
framexbottommargin=0pt,
framexleftmargin=0pt,
framexrightmargin=0pt,
fontadjust=true,
basewidth=0.55em,
upquote=true,
}}

\lstnewenvironment{mypython}[1][]
{
\mypythonstyle
\lstset{#1}
}
{}

\newcommand\mypythonstylenormalinline{\lstset{
language=Python,
basicstyle=\ttfamily\normalsize,
otherkeywords={self},
keywordstyle=\bfseries\ttfamily\color{myteal},
commentstyle=\itshape\color{mymediumteal},
stringstyle=\color{mydeepteal},
emph={MyClass,__init__},
emphstyle=\ttb\color{mydeepteal},
showstringspaces=false,
backgroundcolor=\color{bg},
rulecolor=\color{bg},
breaklines=false,
numbers=left,
numbersep=5pt,
numberstyle=\tiny,
tabsize=4,
xleftmargin=0em,
frame=single,
framesep=3pt,
framextopmargin=0pt,
framexbottommargin=0pt,
framexleftmargin=0pt,
framexrightmargin=0pt,
fontadjust=true,
upquote=true,
}}

\newcommand\mypythoninline[1]{{\mypythonstylenormalinline\lstinline!#1!}}

\lstdefinelanguage{yaml}{
  keywords={true,false,null,y,n},
  keywordstyle=\color{blue}\bfseries,
  basicstyle=\ttfamily\footnotesize,
  comment=[l]{\#},
  morecomment=[s]{/*}{*/},
  commentstyle=\color{gray}\ttfamily,
  stringstyle=\color{red},
  moredelim=[l][\color{orange}]{\&},
  moredelim=[s][\color{orange}]{<}{>},
  morestring=[b]',
  morestring=[b]",
  sensitive=true,
  tabsize=2,
  showspaces=false,
  showstringspaces=false,
  backgroundcolor=\color{codegray},
  frame=single
}
